\chapter{Introduction}\label{ch:intro}

Let $\Lambda$ be the von Mangoldt function and, for a modulus $q \ge 1$ and a residue class
$a \bmod q$, let
$$\psi(x; q, a) := \sum_{\substack{n \le x \\ n \equiv a \pmod q}} \Lambda(n)$$
be Chebyshev's function in the arithmetic progression $a \bmod q$. The prime number theorem in
arithmetic progressions predicts $\psi(x; q, a) \approx x/\varphi(q)$ when $(a, q) = 1$. The best
unconditional estimate for an \emph{individual} modulus, the Siegel--Walfisz theorem,
only covers $q \le (\log x)^C$, and even the generalised Riemann
hypothesis would only reach $q \le x^{1/2 - \varepsilon}$. The theorem of Bombieri
and A.~I.~Vinogradov shows that the latter range is available \emph{on average}
over the modulus: for every $A \ge 0$,
$$\sum_{q \le Q} \max_{y \le x} \max_{a \in (\Z/q\Z)^*} \left| \psi(y; q, a) - \frac{y}{\varphi(q)} \right| \ll_A \frac{x}{(\log x)^A}
\qquad \text{uniformly for } x \ge 2 \text{ and } 1 \le Q \le \frac{\sqrt x}{(\log x)^{A+3}} .$$
This is often described as ``the generalised Riemann hypothesis on average''. It is the main
analytic input in sieve-theoretic applications such as the work of Goldston--Pintz--Y\i ld\i r\i m and
Maynard on small gaps between primes, and it is the reason the level of distribution $1/2$ plays such a
prominent role in that subject.

The original proofs used the large sieve together with zero-density estimates for Dirichlet
$L$-functions. The proof formalised here is the elementary one due to Vaughan,
in the form presented in Chapter~26 of Koukoulopoulos
\cite{Koukoulopoulos2019}.
Vaughan's identity splits $\Lambda$ into three pieces; one is small, one is a \emph{Type~I} sum that
is bounded pointwise by elementary means, and one is a \emph{Type~II} (bilinear) sum that is bounded
on average over $q$ by the large sieve, after the small conductors have been handled with the
Siegel--Walfisz theorem.

\section*{What is formalised}

This blueprint describes a formalisation of the Bombieri--Vinogradov theorem in Lean~4 and Mathlib.
Two analytic results are taken as external inputs, i.e.\ as axioms in Lean: the Siegel--Walfisz
theorem (\Cref{thm:siegel-walfisz}) and the multiplicative large sieve inequality
(\Cref{thm:large-sieve}). Everything else is proved from Mathlib, apart from the smoothed cutoff
function and its basic properties, which are borrowed from the PrimeNumberTheoremAnd library
\cite{PNTplus} (see \Cref{prop:smooth-cutoff}).

\begin{theorem*}[Main result, informal form]
Assume the Siegel--Walfisz theorem and the large sieve inequality. Then for every $A \in \N$, all real
$x \ge 2$ and $1 \le Q \le \sqrt x/(\log x)^{A+3}$,
$$\sum_{q \le Q}\ \sup_{1 \le y \le x}\ \max_{a \in (\Z/q\Z)^*} \left| \psi(y; q, a) - \frac{y}{\varphi(q)} \right| \ll_A \frac{x}{(\log x)^A} .$$
\end{theorem*}
The precise statement is \Cref{thm:bombieri-vinogradov}.

\section*{How to read this document}

The text is organised as a paper, not as a commentary on the Lean code. Each numbered statement
carries exactly the hypotheses that the formal statement carries, and each proof is complete
modulo earlier numbered results; the proofs do, however, follow the route taken by the
formalisation, which occasionally differs from the textbook route (\Cref{sec:deviations}). Below each
statement, the line \emph{Lean} lists the declarations that formalise it. A single mathematical
statement is frequently spread over several Lean declarations, for instance a pointwise bound and its
maximal or summed form, or an identity and the two or three special cases in which it is applied. A
few standard facts that are quoted from Mathlib or from the PrimeNumberTheoremAnd library appear as
numbered results with a one-line proof saying so, so that no proof in this document appeals to an
unnamed ``standard fact''.

Implied constants are not tracked: in Lean every $\ll$ below is witnessed by some real number,
obtained where convenient by a choice or a compactness argument, and nothing is claimed about its
size. In particular the implied constant in the Siegel--Walfisz axiom is not effective, and neither is
any constant that depends on it. \Cref{not:asymptotic} fixes the conventions.
